% concatenated from main.tex
\documentclass[onefignum,onetabnum]{siamart251216}

\usepackage{mathdots}
\usepackage{amsfonts}
\usepackage{amssymb}
\usepackage{latexsym}
\usepackage{graphicx}
\usepackage{pgfplots}
\usepackage{todonotes}
\usepackage{epstopdf}
\usepackage{verbatim}
\usepackage{algorithm}
\usepackage{algorithmic}
\usepackage{hyperref}
\usepackage{amsopn,url}
\usepackage{amsmath}
\usepackage{amstext,amsfonts,bm,amssymb}
\usepackage{indentfirst}
\pgfplotsset{compat=1.18}

\def\subspace{\mathcal{Q}}
\def\Nystrom{Nystr{\"o}m}
\def\svdmethod{SVD-extract}
\def\leadingsubspace{\mathcal{U}_1}
\def\trailingsubspace{\mathcal{U}_2}
\def\projector{\mathcal{P}}
\def\realR{\mathbb{R}}

\newcommand{\eqdef}{:=}
\newcommand{\kernel}{\operatorname{ker}}
\newcommand{\image}{\operatorname{span}}
\DeclareMathOperator{\eig}{eig}
\newcommand{\diag}{{\rm diag}}
\newcommand{\mat}{\mbox{mat}}

\DeclareMathOperator*{\minimize}{minimize}
\DeclareMathOperator*{\maximize}{maximize}
\DeclareMathOperator*{\subjectto}{subject\ to}

\newtheorem{remark}{{\sc Remark}}[section]

\makeatletter
\renewcommand{\@seccntformat}[1]{\csname the#1\endcsname. \,}
\makeatother

\setlength{\marginparwidth}{2cm}



\title{Finding accurate eigenvalues and eigenvectors of positive semi-definite matrices given a subspace}
\author{%
Yuji Nakatsukasa \and
Zheng Tang\thanks{Mathematical Institute, University of Oxford
  (\email{nakatsukasa@maths.ox.ac.uk}, \email{zheng.tang@maths.ox.ac.uk}). YN is supported by the EPSRC grant EP/Y030990/1.}%
}
\headers{Finding accurate eigenvalues given a subspace}{Y. Nakatsukasa and Z. Tang}

\begin{document}

\maketitle
\begin{abstract}
We revisit a classical problem in numerical linear algebra: given an $k$-dimensional subspace $\subspace$ that approximates the leading eigenspace of an $n\times n$ positive semi-definite matrix $A$, the goal is to extract high-accuracy eigenvalues. The Rayleigh-Ritz (RR) method is the standard algorithm for the task, which has been shown to be optimal in several ways (when $A$ is symmetric, not necessarily positive semi-definite $A\succeq 0$). In this paper, we show that when $A\succeq 0$, alternative methods can outperform RR, while having the same computational complexity, that is, the main cost is in computing $AQ$, plus an $O(nk^2)$ term. In particular, we advocate the use of \Nystrom's method, showing that the approximate eigenvalues always have higher accuracy than RR, and the improvement can be arbitrarily large. The difference is significant, especially when $A$ has a fast-decaying spectrum. A similar improvement is numerically observed for the purpose of approximating the leading eigenvectors. In contrast, when the target eigenvalues are the trailing ones, the situation is reversed, and the \Nystrom\ method performs poorly; we suggest a remedy for this situation. 
\end{abstract}

\begin{keywords}
Eigenvalue problem, \Nystrom, positive semi-definite matrix, Rayleigh-Ritz, singular value decomposition, subspace method
\end{keywords}

\begin{MSCcodes}
65F15, 65F10, 15A18
\end{MSCcodes}

\tableofcontents


\section{Introduction}
The large-scale symmetric eigenvalue problem is ubiquitous and computationally challenging in numerical linear algebra. Solving such a problem plays an essential role in a wide range of applications not only in mathematics but also in various engineering fields~\cite{ambikasaran2015fast, axelsson2001finite, boffi2010finite, hu1985toeplitz, mastronardi1999computing, mena2018numerical}, and more recently in the machine learning community for kernel methods~\cite{belabbas2009spectral, gittens2016revisiting, williams2000using}.

In this paper, we mainly focus on the positive semi-definite case $A\succeq 0$. This is an important case for eigenvalue problems in machine learning and scientific computing applications, for example, computing the subset of the eigenpairs of the kernel matrix and the covariance matrix in the kernel methods~\cite{abedsoltan2023toward, fowlkes2004spectral, meanti2020kernel, rudi2017falkon} and principal component analysis~\cite{bishop2006pattern}, respectively. When the matrix is too large for a dense algorithm (e.g., the QR algorithm) to be applicable, the standard approach is to turn to subspace methods. Here, one finds a low-dimensional subspace $\subspace$ (e.g., by a Krylov subspace or the power method and their variants), and extracts eigenvalues and eigenvectors given $\subspace$. Rayleigh-Ritz (RR) has long been the de facto method for the task. In RR, one computes the approximate eigenvalues as those of $Q^\top AQ$, where $Q$ is an $n\times k$ matrix with orthonormal columns spanning the subspace $\subspace$. The complexity is that of computing $AQ$ and then $Q^\top(AQ)$ ($O(nk^2)$), in addition to finding the eigenvalues of $Q^\top AQ$ ($O(k^3)$), so $O(Nk+nk^2)$, where $N$ is the cost of performing a matrix-vector multiplication with $A$. 

When the matrix is positive semi-definite (PSD), additional techniques can be employed to address these problems by exploiting its structural information. Specifically, the PSD property allows us to adopt singular-value-based approaches, as it implies the equivalence between the eigenvalues and the singular values, and more generally, the eigendecomposition and the SVD. In particular, subspace iterations, Krylov subspace methods~\cite[Sec.~12,~14]{parlett1998symmetric}~\cite[Sec.~5,~6]{saad2011numerical}, and randomized techniques~\cite{halko2011finding, martinsson2020randomized} can efficiently approximate the leading singular values of a general matrix, and hence the leading eigenspace of a PSD matrix. 

RR has long been known to be optimal in the sense of the min-max characterization and subspace projection~\cite[Ch.~11]{parlett1998symmetric}. Despite these, in this work we show that there are better techniques to approximate the leading eigenpairs that exploit the positive semi-definiteness, in particular, \Nystrom's method based on the eigenvalues of the \Nystrom\ approximation $\hat{A}_{Nys} = A \langle Q \rangle= AQ (Q^\top A Q)^{\dagger} (AQ)^\top$. The resolution of the apparent inconsistency between the optimality and our result is as follows: the optimality of RR does not imply that the accuracy of the approximate eigenvalues cannot be improved, particularly when the PSD property is used. In this paper, we study the accuracy of the eigenvalue approximations obtained by three methods: RR (based on $\lambda(Q^\top AQ)$), \Nystrom\ (based on $\lambda(AQ (Q^\top A Q)^{\dagger} (AQ)^\top$), and \svdmethod\ (based on $\sigma_i(AQ)$). We explain these methods in more detail in Section~\ref{sec:methods}. The goal of our paper is to theoretically establish that the quality of the approximate eigenvalues is the highest with \Nystrom, followed by \svdmethod\ and RR, in this order. We further show that \Nystrom\ method's accuracy can be better than the other methods by an arbitrarily large factor. The improvement is significant particularly when the spectrum of $A$ decays rapidly. 

We note that a related finding is reported in \cite{lazzarino2025matrix} for the purpose of extracting singular values from approximate (left and right) singular subspaces. This paper can be seen as a specialization to the case $A\succeq 0$, where the simplifications lead to the particularly attractive method of \Nystrom.

Section~\ref{sec:Background-and-Motivation-section} reviews the background information of these techniques for finding the eigenpairs. In Section~\ref{sec:Theoretical-analysis-section}, we study the quality of eigenvalue accuracy extracted by each of the three methods, and show that \Nystrom\ has the highest accuracy, followed by \svdmethod, then RR. Finally, in Section~\ref{sec:Numerical-illustrations-section}, we present numerical experiments illustrating our findings. We also briefly discuss the extraction of trailing eigenvalues, for which \Nystrom\ turns out to be the worst of the three methods. We show that this can be remedied by working with a shift-and-flipped matrix $\gamma I-A$ for $\gamma\geq \lambda_{\max}(A)$, although this requires $\lambda_{\max}(A)$ or its estimate with some extra cost, and the improvement over RR in this case tends to be small.
 
We focus on RR, \svdmethod, and \Nystrom\ in this paper, because they are single-pass methods and have the same computational complexity $O(Nk+nk^2)$. This enables us to avoid addressing the trade-off between method accuracy and complexity. 

\paragraph{Notation}
Unless specified otherwise, for any $m \times n$ matrix $B$, we denote the $i$-th largest singular value by $\sigma_i(B)$ with $\sigma_1(B) \geq \cdots \geq \sigma_{\min\{m,n\}}(B) \geq 0$. Any $n \times n$ PSD matrix $A$ has real nonnegative eigenvalues, which are equal to its singular values. We denote its eigenvalues by $\lambda_{1}(A) \geq \cdots \geq \lambda_{n}(A) \geq 0$ and use $A = U \Lambda U^\top$ to denote its eigendecomposition, where $U =  \left[u_{1}, ..., u_{n}\right]$ is orthogonal and $\Lambda = \diag(\lambda_1, ..., \lambda_n)$. Throughout the paper we use $\leadingsubspace := \image([u_1, ..., u_k])$ and $\trailingsubspace := \image([u_{n-k+1}, ..., u_n])$ to denote the invariant subspace of $A$ spanned by the leading eigenvectors and trailing eigenvectors. We use MATLAB notation for matrix indexing, where $A(i:j, :)$, $A(:, i:j)$ and $A(i:j, i:j)$ denote the matrix consisting the $i$-st to $j$-th columns of $A$, the $i$-st to $j$-th rows of $A$, and their intersection. $\left\| \cdot \right\|_{2}$ denotes the matrix spectral norm and the vector $2$-norm. We say that an $n \times k$ ($n\geq k$) matrix $Q$ is orthonormal if it has orthonormal columns such that $Q^\top Q = I_{k}$. The orthogonal complement of an orthonormal matrix $Q$ is denoted by $Q_{\perp}\in\mathbb{R}^{n\times (n-k)}$. We use $B^{\dagger}$ to denote the pseudoinverse of any $m \times n$ matrix $B$, which is defined via the 
SVD $B=U_*\Sigma_*V_*^\top = \sum_{i=1}^{r}\sigma_{i}(B) u_{*,i}  v_{*,i}^\top $, where $U_{*} = \left[u_{*,1}, ..., u_{*,n}\right]$, $V_{*} = \left[v_{*,1}, ..., v_{*,n}\right]$ and $r=\mbox{rank}(B)$, as 
\[
    B = \sum_{i=1}^{r}\sigma_{i}(B)  u_{*,i}  v_{*,i}^\top \quad \text{ then }\quad B^{\dagger} = V_{*} \diag (\sigma_{1}^{-1}, ..., \sigma_{r}^{-1}, 0, ..., 0) U_{*}^\top,
\]
where we use $*$ to avoid confusion between the left singular vector and eigenvector. The symbol $\succeq$ denotes the Loewner order of matrices: $A \succeq B$ means $A-B$ is positive semi-definite. $A\langle X \rangle$ represents the \Nystrom\ approximation of a PSD matrix $A$ with respect to a test matrix $X\in\mathbb{R}^{n\times k}$ such that $A\langle X \rangle = AX (X^\top A X)^{\dagger} X^\top A$. For simplicity, we work with real matrices; however, analogous results hold for complex matrices, possibly differing slightly in constant terms.


\section{Background}
\label{sec:Background-and-Motivation-section}
Subspace methods first approximate an invariant subspace and then extract the eigenpairs given that approximation. Our core concern is the second step and the goal of this paper is to exploit the matrix's PSD property to improve the accuracy of eigenpair extraction beyond RR. This section provides an introduction and offers a first look at the comparison of these methods.

\subsection{Projection-based methods to extract eigenpairs given a subspace}\label{sec:methods}
Here, we provide an overview of three projection-based approaches, including RR, to accomplish the extraction task.

\subsubsection{Rayleigh-Ritz (RR)}
RR, described in Algorithm~\ref{alg:Rayleigh-Ritz}, is classic and the most widely method based on projecting the matrix and computing the eigenvalue decomposition of $Q^\top A Q$.

\begin{algorithm}
	\algsetup{linenosize=\small} \small
	\renewcommand{\algorithmicrequire}{\textbf{Input:}}
	\renewcommand{\algorithmicensure}{\textbf{Output:}}
	\caption{RR}
	\begin{algorithmic}[1]
		\label{alg:Rayleigh-Ritz}
		\REQUIRE Positive semi-definite $A \in \mathbb{R}^{n \times n}$; orthonormal $Q \in \mathbb{R}^{n \times k}$.
		\ENSURE Diagonal $\hat{\Lambda} \in \mathbb{R}^{k \times k }$ and orthonormal $\hat{U} \in \mathbb{R}^{n \times k}$ that approximate the desired eigenpairs, i.e., $A\hat{U} \approx \hat{U}\hat{\Lambda}$
		\STATE Compute $Q^\top A Q$.
		\STATE Compute the eigendecomposition $Q^\top A Q = \widetilde{\Omega} \hat{\Lambda} \widetilde{\Omega}^\top$, where $\tilde{\Omega}$ is orthogonal.
		\STATE $\hat{U} = Q \widetilde{\Omega}$.
	\end{algorithmic} 
\end{algorithm}


\subsubsection{\svdmethod}
The positive semi-definiteness implies the equivalence between eigenpairs and singular pairs. In (randomized) numerical linear algebra, the one-side projected SVD (we abbreviate it as \svdmethod) technique is a common way to approximate the singular pairs (so eigenpairs) of a matrix $A$ when a subspace $\subspace$ is given~\cite{gu2015subspace, halko2011finding, martinsson2020randomized}. This allows us to obtain eigenvalue approximations for $A$ by computing the singular values of the one-side projected matrix $AQ$ (or $AQQ^\top$)
\[
    \hat{\lambda}_i = \sigma_i(AQ).
\]
The approximations to eigenvectors have two variants: once the SVD
\[
    AQ = \widetilde{U}_{*} \widetilde{\Sigma}_* \widetilde{V}_{*}^\top
\]
is computed, the approximation can be either the columns of $Q\widetilde{V}_{*}$ or $\widetilde{U}_{*}$. The former extracts eigenvectors from $\image(Q)$ as approximations, as does the RR method, whereas the latter seeks the eigenvector approximations from the subspace $\image(AQ)$. We distinguish these two variants of \svdmethod\ for singular vector approximation by denoting them as SVD-$Q\tilde{V}_{*}$ and SVD-$\tilde{U}_{*}$. Algorithm~\ref{alg:SVD} provides the pseudocode of the \svdmethod. 

\begin{algorithm}
	\algsetup{linenosize=\small} \small
	\renewcommand{\algorithmicrequire}{\textbf{Input:}}
	\renewcommand{\algorithmicensure}{\textbf{Output:}}
	\caption{\svdmethod}
	\begin{algorithmic}[1]
		\label{alg:SVD}
		\REQUIRE Positive semi-definite $A \in \mathbb{R}^{n \times n}$; orthonormal $Q \in \mathbb{R}^{n \times k}$.
		\ENSURE Diagonal $\hat{\Lambda} \in \mathbb{R}^{k \times k }$ and orthonormal $\hat{U} \in \mathbb{R}^{n \times k}$ that approximate the desired eigenpairs, i.e., $A\hat{U} \approx \hat{U}\hat{\Lambda}$
		\STATE Compute $AQ$.
		\STATE Compute the SVD $AQ = \widetilde{U}_{*} \widetilde{\Sigma}_* \widetilde{V}_{*}^\top$.
		\STATE $\hat{\Lambda} = \widetilde{\Sigma}_*$, and either $\hat{U} = Q \widetilde{V}_{*}$ or $\hat{U} = \widetilde{U}_{*}$.
	\end{algorithmic} 
\end{algorithm}

\subsubsection{\Nystrom}
Originally introduced in machine learning~\cite{williams2000using}, and now a popular tool in randomized numerical linear algebra, \Nystrom\ is an attractive method to find a low-rank approximation for a PSD matrix and thus its eigenpairs \cite{halko2011finding, martinsson2020randomized, nakatsukasa2020sharp, tropp2023randomized}. It seeks eigenvalue approximations by
\[
    \hat{\lambda}_i = \lambda_i \left(A \langle Q \rangle\right) = \lambda_i\left(AQ (Q^\top A Q)^{\dagger} (AQ)^\top\right)
\]
and the corresponding eigenvectors of $A \langle Q \rangle$ are eigenvector approximations. The pseudocode of the \Nystrom\ method appears in Algorithm~\ref{alg:Nystrom}. For efficiency, it implicitly generates the \Nystrom\ approximation by $(AQL^{\dagger})(AQL^{\dagger})^\top$ with the Cholesky factorization $Q^\top A Q = L^\top L$. The eigenpair approximations are obtained by computing the singular pairs of $\widetilde{R}(Q^\top A Q)^{\dagger}\widetilde{R}^\top$, where $AQ = \widetilde{Q} \widetilde{R}$ is a thin QR factorization. Specifically, we solve $ZL = \widetilde{R}$ \footnote{We will never need to explicitly form the pseudoinverse, which is computationally expensive when the matrix size is large.} to implicitly generate $\tilde{R} L^{\dagger}$, factorize
\[
    \widetilde{R} C^{\dagger} = \widehat{U}_{*} \widehat{\Sigma}_* \widehat{V}_{*}^\top,
\]
and obtain eigenpair approximations 
\[
    \hat{\Lambda} = \widehat{\Sigma}^2, \hat{U} = \widetilde{Q} \widehat{U}_{*}.
\]
In finite-precision computation, the implementation of \Nystrom\ may suffer from numerical instability. This is due to the ill-conditioning of the core matrix $Q^\top A Q$, which makes the computation of its pseudoinverse highly sensitive to round-off errors. To our knowledge, there are two main classes of proposed techniques to overcome this instability. One is to introduce an $\epsilon$-truncated Cholesky factorization for computing the pseudoinverse \cite{bucci2025numerical}, which is shown to be stable when $S$ is a subsampling matrix that gives a local-maximum volume submatrix~\cite{damle2025estimating}, and observed to work well in practice more generally.  Another approach~\cite{tropp2017fixed} is to incorporate a small shift to ensure the smallest singular value of the core matrix is larger than the unit round-off. In our paper we choose the implementation of the former one, as it is more efficient and stability was not seen to cause issues (even without the $\epsilon$ truncation).

\begin{remark}
    The \Nystrom\ approximation $A \langle Q \rangle$ only depends on the range of the sketching matrix $Q$. That is, for any non-singular $M\in\mathbb{R}^{k\times k}$, we have
    \[
        AQ (Q^\top AQ)^{\dagger} (AQ)^\top = A(QM) \left((QM)^\top A (QM)\right)^{\dagger} \left(A (QM)\right)^\top.
    \]
\end{remark}

\begin{algorithm}
    \algsetup{linenosize=\small} \small
    \renewcommand{\algorithmicrequire}{\textbf{Input:}}
    \renewcommand{\algorithmicensure}{\textbf{Output:}}
    \caption{\Nystrom\ method}
    \begin{algorithmic}[1]
    \label{alg:Nystrom}
    \REQUIRE Positive semi-definite $A \in \mathbb{R}^{n \times n}$; orthonormal $Q \in \mathbb{R}^{n \times k}$; a truncation parameter $\epsilon > 0$
    \ENSURE Diagonal $\hat{\Lambda} \in \mathbb{R}^{k \times k }$ and orthonormal $\hat{U} \in \mathbb{R}^{n \times k}$ that approximate the desired eigenpairs, i.e., $A\hat{U} \approx \hat{U}\hat{\Lambda}$
    \STATE Compute $Y = AQ$ and its QR factorization $Y = \widetilde{Q}\widetilde{R}$.
    \STATE Compute $W = Q^\top Y$.
    \STATE Compute $\epsilon$-truncated Cholesky factorization $W = L_{\epsilon}^\top L_{\epsilon}$ (in practice, standard Cholesky $W = L^\top L$ suffices).
    \STATE Solve the overdetermined linear systems $Z L_{\epsilon} = \widetilde{R}$ for $Z$ via back substitution.
    \STATE Compute the SVD to $Z = \widehat{U}_{*} \widehat{\Sigma}_* \widehat{V}_{*}^\top$.
    \STATE $\hat{\Lambda} = \widehat{\Sigma}^2$, $\hat{U} = \widetilde{Q}\widehat{U}_{*}$.
    \end{algorithmic} 
\end{algorithm}

\subsection{Connection among RR, \svdmethod and \Nystrom}
\label{sec:Compare-approximate-eigenvalue-RR-vs-SVD-vs-Nystrom}
The main difference among these methods lies in the subspace where they extract the eigenvectors as approximations. RR and SVD-$Q\widetilde{V}$ extract the eigenvectors from the given subspace $\image(Q)$, whereas the SVD-$\widetilde{U}_{*}$ and \Nystrom\ seek eigenvectors from $\image(AQ)$. Let us show that the RR method can be viewed as an \svdmethod\ method applied to a matrix other than the original one and vice versa. Define the square root of a PSD matrix $A = U \Sigma U^\top$ via the SVD
\[
    A^{1/2} := U \Lambda^{1/2} U^\top,
\]
where $\Lambda^{1/2} = \text{diag} (\lambda_1(A)^{1/2}, ..., \lambda_n(A)^{1/2})$. From the \svdmethod\ perspective, the RR eigenvalue approximations for $A$ can be viewed as
\[
    \hat{\lambda}_{i}^{(\text{RR})}(A) = \lambda_{i}(Q^\top A Q) = \left(\sigma_{i}(A^{1/2}Q)\right)^2 = \left(\hat{\sigma}_{i}^{(\text{SVD})}(A^{1/2})\right)^2,
\]
where $i = 1, ..., \dim(\image(Q))$. That is, the RR eigenvalue estimates are equal to those of (the square of) SVD-extract applied to $A^{1/2}$ with the same $Q$. Similarly, from the RR perspective, the \svdmethod\ approximation to the singular values of $A$ can be viewed as the square root of RR eigenvalue approximations for $A^2$, that is,
\[
    \hat{\sigma}_{i}^{(\text{SVD})}(A) = \sigma_{i}(AQ) = \sqrt{\lambda_{i}(Q^\top A^2 Q)} = \sqrt{\hat{\lambda}_{i}^{(\text{RR})}(A^2)},
\]
where $i = 1, ..., \dim(\image(Q))$. 

Also, when $\subspace$ is an approximation to the invariant subspace $\leadingsubspace$ of $A$, we can view the RR and \Nystrom\ methods as seeking the eigenpairs of different projection-based sketches to $A$ as approximations. In detail, the RR and \Nystrom\ methods seek the leading eigenpairs of
\[
    \hat{A}_{\text{RR}}^{(k)} = Q Q^\top A Q^\top Q = \left(\mathcal{P}_{\image(Q)} A^{1/2}\right) \left(A^{1/2} \mathcal{P}_{\image(Q)}\right)
\]
and
\[
    \hat{A}_{\text{Nys}}^{(k)} = (AQ) (Q^\top A Q)^{\dagger} (AQ)^\top = \left(A^{1/2} \mathcal{P}_{\image(A^{1/2}Q)}\right) \left(\mathcal{P}_{\image(A^{1/2}Q)} A^{1/2}\right)
\]
as approximations, where $\mathcal{P}_{\image(Q)}$ and $\mathcal{P}_{\image(A^{1/2}Q)}$ is the orthogonal projector onto the subspace $\image(Q)$ and $\image(A^{1/2}Q)$, respectively \cite{gittens2016revisiting}. The latter subspace benefits from one extra step of the block power method, so it more accurately captures the invariant subspace $\leadingsubspace$ than the former; we make this precise in Section~\ref{subsec:Eigenvector-analysis-subsection}.


\subsection{Computational cost}
Computational cost is an important aspect of performance. Here, we measure the cost in terms of the floating-point operations. Table~\ref{table:method-complexity-table} summarizes the complexity of these three methods, including the constants. All methods are $O(Nk+nk^2)$, thereby allowing us to largely avoid the trade-off between method accuracy and complexity. 

\begin{table}[ht]
    \centering
    \caption{The complexity of different methods, where $k = \dim(\subspace)$ and $N$ is the cost of performing a matrix-vector multiplication with $A$.
    }
    \label{table:method-complexity-table}
    \begin{tabular}{c|c}
        \hline
        & Method complexity \\
        \hline
        Rayleigh-Ritz & $Nk + 2nk^2$ \\
        \svdmethod & $Nk + 4nk^2$ \\
        \Nystrom & $Nk+ 4nk^2$ \\
        \hline
    \end{tabular}
\end{table}

While the cost of RR is slightly lower than \Nystrom, we will see that the latter can give significantly better approximations to the eigenpairs, making it well worth the extra cost. 


\section{Theoretical analysis}
\label{sec:Theoretical-analysis-section}
In this section, we will show that  \Nystrom\ has the highest accuracy for approximating the leading eigenvalues, while the RR method is the best for approximating the trailing ones. We will also demonstrate the higher-order accuracy of the \Nystrom\ method in leading eigenvalue approximation, where this improvement can be significant when the spectrum of $A$ decays rapidly.

\subsection{\Nystrom\ works better for leading eigenvalue approximation}
\label{subsec:Leading-eigenvalue-subsection}
We first provide a qualitative comparison between the three methods.
\begin{theorem}
    \label{thm:Nystom-vs-SVD-vs-RR-Theorem}
    Suppose $A$ is an $n \times n$ PSD matrix and $Q\in\mathbb{R}^{n\times k}$ is an orthonormal basis of the subspace $\subspace$ which approximates the invariant subspace $\leadingsubspace$ of $A$. Then, the eigenvalue approximations from the RR, \svdmethod, and \Nystrom\ methods satisfy  
    \begin{equation}
        \label{eq:Nystrom-vs-SVD-vs-RR}
        \lambda_i(A) \geq \lambda_i(A\langle Q \rangle) \geq \sigma_i(AQ) \geq \lambda_i(Q^\top A Q) \quad \forall i = 1,\ldots, k.
    \end{equation}
\end{theorem}

\begin{proof}
    The first inequality follows immediately from $A - A\langle Q \rangle \succeq 0$, which is a consequence of the fact that the Schur complement of an PSD matrix is PSD (more details can be found for example in \cite[Section 4.6]{tropp2023randomized}). To prove the inequality between the RR and \svdmethod\, we have
    \[
        \lambda_{i} (Q^\top A Q) = \sigma_{i} (Q^\top A Q) \leq \sigma_{i} (AQ) \left\|Q^\top\right\|_2 = \sigma_{i} (AQ).
    \]
    
    We now prove the inequality between the \Nystrom\ and \svdmethod. Since $\subspace$ is an $\dim(\subspace)$-dimensional subspace of $\mathbb{R}^{n}$, by applying an orthogonal change of basis with respect to $Q_F := [Q,Q_\perp]$, we can assume
    \[
        \image(Q) =
        \image\left(
        \begin{bmatrix}
            I_{\dim(\subspace)} \\
            0
        \end{bmatrix}
        \right)
        \Rightarrow
        Q = 
        \begin{bmatrix}
            I_{\dim(\subspace)} \\
            0
        \end{bmatrix}
        M,
    \]
    where $M$ is $k\times k$ orthogonal and we denote $X := \begin{bmatrix}
        I_{\dim(\subspace)} & 0
    \end{bmatrix}^\top$ for convenience. Since $A \langle Q \rangle$ only depends on the range of $Q$, we have 
    \[
        \begin{aligned}
            & \lambda_{i} (Q^\top AQ) = \lambda_{i} (M^\top X^\top A X M) = \lambda_{i} (X^\top A X^\top), \\
            & \sigma_{i} (AQ) = \sigma_{i} (AXM) = \sigma_{i} (AX), \\
            & \lambda_{i} \left(A \langle Q \rangle\right) = \lambda_{i} \left(A(XM) \left((XM)^\top A (XM)\right)^{\dagger} (XM)^\top A^\top\right)  = \lambda_{i} (A \langle X \rangle).
        \end{aligned}
    \]
    Then, we can assume without loss of generality that
    \[
        Q = X
        =
        \begin{bmatrix}
        I_{\dim(\subspace)} \\
        0 \\
        \end{bmatrix}.
    \]
    To prove $\lambda_i(A\langle Q \rangle) \geq \lambda_i(A\langle Q \rangle)$, we consider the matrices in block form:
    \[
    \begin{gathered}
        A =
        \begin{bmatrix}
            A_{11} & A_{12} \\
            A_{12}^\top & A_{22}
        \end{bmatrix},
        \qquad
        AX =
        \begin{bmatrix}
            A_{11} \\
            A_{12}^\top
        \end{bmatrix},
        \\[0.5em]
        A \langle X \rangle
        =
        AX (X^\top A X)^{\dagger} X^\top A^\top
        =
        \begin{bmatrix}
            A_{11} & A_{12} \\
            A_{12}^\top & A_{12}^\top A_{11}^{\dagger} A_{12}
        \end{bmatrix}.
    \end{gathered}
    \]
    Thus, we have
    \[
        \begin{aligned}
            \lambda_i(A\langle X \rangle)
            =
            \lambda_{i} \left(
                        \begin{bmatrix}
                        A_{11} & A_{12} \\
                        A_{12}^\top & A_{12}^\top A_{11}^{\dagger} A_{12} \\
                        \end{bmatrix}
                        \right)
            & =
            \sigma_{i} \left(
                        \begin{bmatrix}
                        A_{11} & A_{12} \\
                        A_{12}^\top & A_{12}^\top A_{11}^{\dagger} A_{12} \\
                        \end{bmatrix}
                        \right)
            \\
            & \geq
            \sigma_{i} \left(
                       \begin{bmatrix}
                       A_{11} \\
                       A_{12}^\top \\
                       \end{bmatrix}
                       \right)
            =
            \sigma_i(AX).
        \end{aligned}
    \]
    The second equality is due to $A \langle X \rangle \succeq 0$. This completes the proof.
\end{proof}

\begin{remark}
    The inequalities between the eigenvalue approximations do not require $\subspace$ to be an invariant subspace approximation, i.e., in the same settings as Theorem~\ref{thm:Nystom-vs-SVD-vs-RR-Theorem} except that $\mathcal{Q}$ is a low-dimensional subspace, we have
    \begin{equation}
        \label{remark:Nystrom-vs-SVD-vs-RR-general-form}
        \lambda_i(A\langle Q \rangle) \geq \sigma_i(AQ) \geq \lambda_i(Q^\top A Q) \quad \forall i = 1, ... ,\dim(\subspace).
    \end{equation}
\end{remark}
Theorem~\ref{thm:Nystom-vs-SVD-vs-RR-Theorem} shows that the leading eigenvalue approximations
\[
\hat{\lambda}_{i}^{\mathrm{RR}}
    := \lambda_i(Q^\top A Q), 
\qquad
\hat{\lambda}_{i}^{\mathrm{SVD}}
    := \sigma_i(AQ),
\qquad
\hat{\lambda}_{i}^{\mathrm{Nys}}
    := \lambda_i(A\langle Q\rangle)
\]
satisfy the following bounds.
\begin{equation}
    \label{eq:largest-eigenvalue-accuracy-relation}
    \left| \lambda_{i}(A) - \hat{\lambda}_{i}^{\text{(RR)}}\right| 
    \geq
    \left| \lambda_{i}(A) - \hat{\lambda}_{i}^{\text{(SVD)}}\right|
    \geq
    \left| \lambda_{i}(A) - \hat{\lambda}_{i}^{\text{(Nys)}}\right|,
\end{equation}
where $i = 1, ..., \dim(\leadingsubspace)$. This implies that exploiting the PSD property does help the \svdmethod\ and \Nystrom\ improve the accuracy of approximating the leading eigenvalues.


\subsection{High-order accuracy of \Nystrom\ in leading eigenvalue approximations}
\label{subsec:Analysis-High-order-accuracy-Nystrom}
The results above show qualitatively that \Nystrom\ and \svdmethod\ are better than RR in approximating the leading eigenvalues. Here we derive a quantitative result, showing that \Nystrom\ has a higher-order accuracy than \svdmethod\, and RR, and this accuracy improvement can be arbitrarily large, especially in the leading eigenvalue approximations.

\begin{theorem}
    \label{thm:Higher-order-accuracy-Nystrom-theorem}
    Suppose that $A$ is an $n \times n$ PSD matrix and that $\subspace$ is an $k$-dimensional $\epsilon$-approximation, with $\epsilon < 1$, to the dominant invariant subspace $\leadingsubspace$ of $A$. Let $Q$ be an orthonormal basis matrix for $\subspace$ such that\footnote{$\subspace$ is an $\epsilon$-approximation to $\leadingsubspace$ in the principal-angle sense
    \[
    \|\sin \Theta(\image(U_1), \image(Q))\|_2 = \|U_2^\top Q\|_2 = \epsilon \|N\|_2 \leq \epsilon,
    \]
    which requires $\|N\|_2 \leq 1$. $\|M\|_2$ controls basis alignment inside $\image(U_1)$ since
    \[
    U_1^\top Q = I_k + \epsilon^2 M.
    \]
    If we restrict $\|M\|_2 = O(1)$, the basis $Q$ is aligned with $U_1$ up to a second-order perturbation inside $\image(U_1)$ such that
    \[
    \|U_1^\top Q - I_k\|_2 = O(\epsilon^2).
    \]
    It is convenient to normalize it by $\|M\|_2 \leq 1$, which gives the explicit estimate
    \[
    \|U_1^\top Q - I_k\|_2 \leq \epsilon^2.
    \]
    }
    \begin{equation}
        \label{eq:eps-approximation-Q-formula}
        Q = U_1 + \epsilon^2 U_1 M + \epsilon U_2 N,
    \end{equation}
    where
    \[
        U_1 = [u_1,\ldots,u_k], 
        \qquad
        U_2 = [u_{k+1},\ldots,u_n],
    \]
    and
    \[
        M \in \realR^{k \times k}, 
        \qquad
        N \in \realR^{(n-k) \times k},
        \qquad
        \|M\|_2 \leq 1, 
        \qquad
        \|N\|_2 \leq 1 .
    \]

    Define 
        \[
\alpha_i :=   1-\frac{  3\lambda_{k+1}(A) + 3 \epsilon^2 (\lambda_1(A) - \lambda_{k+1}(A))}{\lambda_i(A)} . 
    \]
    For each $i=1,\ldots, k$, if 
    $\alpha_i>0$, we have
    \begin{equation}
        \label{eq:Upper-bound-RR}
        \left|\lambda_{i}(A) - \sigma_i(AQ)\right| \leq \left|\lambda_{i}(A) - \lambda_i(Q^\top A Q)\right| \leq C_{RR} \cdot \epsilon^2 \lambda_{1}(A),
    \end{equation}
    \begin{equation}
        \label{eq:Upper-bound-Nys}
        \left| \lambda_{i}(A) - \lambda_{i}(A \langle Q \rangle) \right| \leq C_{Nys,i} \cdot \left(\epsilon^2\lambda_{k+1}(A) + \epsilon^4 (\lambda_1(A) - \lambda_{k+1}(A))\right),
    \end{equation}
    where 
    \[
        C_{RR} := 2 + \epsilon^2 + \lambda_{k+1}(A) / \lambda_1(A) \leq 4
    \]
    and
    \[
        C_{Nys, i} := \left(\lambda_1(A) (1+ \epsilon^2) + \lambda_{k+1}(A) \left(1 + \epsilon^2 \left(1 / \left(1 + \sqrt{1-\epsilon^2}\right)\right)\right) / \alpha_{i} \lambda_i(A)\right)^2.
    \]
\end{theorem}

\begin{proof}
Let us assume without losing generality that the PSD matrix $A$ is diagonal such that
\[
    A
    =
    \diag (\lambda_{1}, ..., \lambda_{n})
    =
    \begin{bmatrix}
        \Lambda_{1} & 0 \\
        0 & \Lambda_{2} \\
    \end{bmatrix},
\]
where $\Lambda_{1} = \diag(\lambda_1, ..., \lambda_k)$ is an $k \times k$ diagonal matrix. Unless otherwise specified, in the following proofs, we use $\lambda_i$ to denote the eigenvalues $\lambda_i(A)$ of $A$ and $\lambda_i(B)$ to denote the eigenvalues of any matrix $B$ other than $A$. Since the orthogonal basis $Q$ of $\subspace$ is an $\epsilon$ approximation to the leading invariant subspace of $A$, under the orthogonal change of basis we have \footnote{The first block of $Q$ after changing of basis should be taken to be $I + O(\epsilon^2)$ instead of $I + O(\epsilon)$, since we need $Q^\top Q = I$. More precisely: it could be $I+S$ for a skew-symmetric matrix $S$ with $\|S\|=O(\epsilon)$; one can then right-multiply a $k\times k$ orthogonal matrix $I-S+O(\epsilon^2)$ to obtain another basis $Q$ whose first block is $I + O(\epsilon^2)$. Finally, recall that the output of RR, \svdmethod, and \Nystrom\ are the same as long as the span of $Q$ is the same.
}
\begin{equation}
    \label{eq:transformed-Q-block-form}
    Q = U^T (U_1 + \epsilon^2 U_1 M + \epsilon U_2 N) =
    \begin{bmatrix}
        I_k + \epsilon^2 M \\
        \epsilon N \\
    \end{bmatrix}.
\end{equation}
In this case, we choose the basis of the orthogonal complement $Q_{\perp}$ of $Q$ so that it is in the form of
\[
    Q_{\perp} =
    \begin{bmatrix}
        \epsilon \widetilde{M} \\
        I_{n-k} + \epsilon^2 \widetilde{N}
    \end{bmatrix}
\]
with $\|\widetilde{M}\|_2 \leq 1$ and $\|\widetilde{N}\|_2 \leq 1 / (1 + \sqrt{1 - \epsilon^2})$ \footnote{Here we use the Cosine-Sine decomposition \cite{stewart1977perturbation} to fix the orthogonal complement and derive the upper bound for $\|\widetilde{N}\|_2$. There exists orthogonal matrices $W \in \realR^{k \times k}, Z \in \realR^{(n-k) \times (n-k)}, Y \in \realR^{k \times k}$ and diagonal matrices $\cos \Theta = \diag(\cos \theta_1, ..., \cos \theta_k), \sin \Theta = \diag(\sin \theta_1, ..., \sin \theta_k) \in \realR^{k \times k}$ with $\cos^2 \Theta + \sin^2 \Theta = I_k$ such that
\[
    \begin{bmatrix}
        Q & Q_{\perp}
    \end{bmatrix}
    =
    \begin{bmatrix}
        W & 0 \\
        0 & Z
    \end{bmatrix}
    \begin{bmatrix}
        \cos \Theta & -\sin \Theta & 0 \\
        \sin \Theta & \cos \Theta & 0 \\
        0 & 0 & I_{n-2k}
    \end{bmatrix}
    \begin{bmatrix}
        Y & 0 \\
        0 & I_{n-k}
    \end{bmatrix}.
\]
Since $\sin \theta_i \leq \max_i \sin \theta_i = \left\|
\begin{bmatrix}
    \sin \Theta \\
    0
\end{bmatrix}
\right\|_2 = \|\epsilon N\|_2 \leq \epsilon$, we have $1 - \cos \theta_i \leq 1 - \sqrt{1-\epsilon^2} = \epsilon^2 / (1 + \sqrt{1-\epsilon^2})$ and further $\|I - (I + \epsilon^2 \widetilde{N})\|_2 = \max\{\max_i|1 - \cos \theta_i|, 0\} = \max_i (1 - \cos \theta_i) \leq \epsilon^2 / (1 + \sqrt{1-\epsilon^2})$.
}.

By Weyl's inequality, the RR approximation to the leading eigenvalues is $O(\epsilon^2 \lambda_1)$ such that
\[
    \begin{aligned}
        \left|\lambda_{i}(A) - \lambda_{i}\left(Q^\top A Q\right)\right| & \leq \|\Lambda_1 - Q^\top A Q\|_2 \\ 
        & = \|-\epsilon^2 \left(M^\top \Lambda_1 + \Lambda_1 M - N^\top \Lambda_2 N\right) - \epsilon^4 M^\top \Lambda_1 M\|_2 \\
        & \leq 2 \epsilon^2 \lambda_1 + \epsilon^4 \lambda_1 + \epsilon^2 \lambda_{k+1} 
        \\
        & = \epsilon^2 \lambda_1 \left(2 + \epsilon^2 + \frac{\lambda_{k+1}}{\lambda_1}\right) \\
        & = C_{RR} \cdot \epsilon^2 \lambda_1, \quad i=1,...,k,
    \end{aligned}
\]
where the constant depends on $A$ such that $C_{RR} := 2 + \epsilon^2 + \lambda_{k+1} / \lambda_1 \leq 4$.
Then, by Theorem~\ref{thm:Nystom-vs-SVD-vs-RR-Theorem} we have
\[
    \left|\lambda_{i}(A) - \sigma_{i}\left(A Q\right)\right| \leq \left|\lambda_{i}(A) - \lambda_{i}\left(Q^\top A Q\right)\right| \leq C_{RR} \cdot \epsilon^2 \lambda_1, \quad i=1,...,k.
\]

Finally, we analyze the accuracy of the \Nystrom\ eigenvalue approximations. Define an $n \times n$ orthogonal matrix $Q_F$ by $Q_{F} = \begin{bmatrix} Q & Q_{\perp} \\ \end{bmatrix}$. Under similarity transformation, the \Nystrom\ approximation error $A - A \langle Q \rangle$ has non-zero structure in the $Q_{F}$ coordinate such that
\begin{equation}
    \label{eq:eigenvalue-perturbation-eqn}
    \begin{aligned}
        & Q_{F}^\top (A - A \langle Q \rangle) Q_{F} \\
        & =
        \begin{bmatrix}
        Q^\top A Q & Q^\top A Q_{\perp} \\
        Q_{\perp}^\top A Q & Q_{\perp}^\top A Q_{\perp}
        \end{bmatrix}
        \\
        & -
        \begin{bmatrix}
        Q^\top A Q & (Q^\top A Q) (Q^\top A Q)^{\dagger} (Q^\top A Q_{\perp}) \\
        (Q_{\perp}^\top A Q) (Q^\top A Q)^{\dagger} (Q^\top A Q) & (Q_{\perp}^\top A Q) (Q^\top A Q)^{\dagger} (Q^\top A Q)
        \end{bmatrix} \\
        & =
        \begin{bmatrix}
            0 & 0 \\
            0 & F
        \end{bmatrix},
    \end{aligned}
\end{equation}
where $F := Q_{F}^\top A Q_{F}/ Q^\top A Q$ is the Schur complement of $Q^\top A Q$ in $Q_{F}^\top A Q_{F}$ and we make the generic assumption that $Q^\top A Q$ is full rank \footnote{So, the pseudoinverse is simply the inverse $(Q^\top A Q)^{\dagger}=(Q^\top A Q)^{-1}$; this assumption holds as long as $Q$ is orthonormal in the form of \eqref{eq:transformed-Q-block-form} and either 
\begin{enumerate}
    \item $A \succ 0$; or 
    \item $A \succeq 0$ but $\lambda_k > 0$ with $\epsilon < 1$; this yields $\lambda_{\min} (Q^\top A Q) \geq \lambda_{min} \left((I + \epsilon^2 M)^\top \Lambda_1 (I + \epsilon^2 M)\right) \geq \lambda_{k} \|I + \epsilon^2 M\|_2^2 = \lambda_{k} (1 - \|\epsilon N\|_2^2) \geq \lambda_{k} (1 - \epsilon^2) > 0$.
\end{enumerate}}. In other words, the differences between the leading $k$ eigenvalues of $A$ (exact values) and those of $A \langle Q \rangle$ (approximate values) only results from the perturbation $F$ to
\[
    Q_{F}^\top A Q_{F}
    =
    \begin{bmatrix}
        Q^\top A Q & E^\top \\
        E & Q_{\perp}^\top A Q_{\perp} \\
    \end{bmatrix},
\]
where
\[
    \|E\|_2 = \|Q_{\perp}^\top A Q\|_2 \leq \epsilon \lambda_1 \left(1 + \epsilon^2\right) + \epsilon \lambda_{k+1} \left(1 + \epsilon^2 \left(\frac{1}{1 + \sqrt{1 - \epsilon^2}}\right)\right).
\]
Since $Q_{F}^\top A Q_{F}$ is a nearly block diagonal matrix, the difference between the leading $i$-th eigenvalue of the block Hermitian matrix $Q_{F}^\top A Q_{F}$ and the matrix after it is perturbed by \eqref{eq:eigenvalue-perturbation-eqn} (i.e. $Q_{F}^\top A \langle Q \rangle Q_{F}$) is bounded above by $O\left( (\left\|E\right\|_2 / \lambda_{i})^2 \left\|F\right\|_2\right)$, as is studied in \cite[Theorem 3.2]{nakatsukasa2012eigenvalue}. More precisely,
\begin{equation}
    \label{eq:Nystrom-perturbation-eq}
    \left| \lambda_{i}(Q_{F}^\top A Q_{F}) - \lambda_{i}(Q_{F}^\top A \langle Q \rangle Q_{F}) \right|
    \leq
    \|F\|_2 \left(\frac{\|E\|_2}{\min|\lambda_i - \lambda(Q_{\perp}^\top A Q_{\perp})| - 2 \|F\|_2}\right)^2.
\end{equation}
That is, the leading eigenvalues of $Q_{F}^\top A Q_{F}$ are extremely insensitive to the structured perturbation \eqref{eq:eigenvalue-perturbation-eqn}. 

We now only need an upper bound for $\left\|F\right\|_2$ to bound \eqref{eq:Nystrom-perturbation-eq}. Let 
\[
    Q_F^\top A Q_F = L L^\top
\]
be the full Cholesky factorization of the PSD matrix $Q_F^\top A Q_F$, then the perturbation $F$ is equivalent to the matrix outer product \cite{zhang2006schur}: using MATLAB notation, 
\[
    L / L(1:k, 1:k) \cdot (L / L(1:k, 1:k) )^\top = L(k:n, k:n) \cdot L(k:n, k:n)^\top.
\]
In this case, $F$ can be seen as the remainder term after $k$ steps of applying the pivoted partial Cholesky factorization to $Q_{F}^\top A Q_{F}$ \cite{chen2025randomly}. Then, we have
\[
    F \succeq 0, \quad Q_{\perp}^\top A Q_{\perp} - F = Q_{\perp}^\top A \langle Q \rangle Q_{\perp} \succeq 0
\]
and further
\[
    \begin{aligned}
        \left\|F\right\|_2
        & \leq
        \left\|Q_{\perp}^\top A Q_{\perp}\right\|_2 \\
        & =
        \max_{x \in \image(Q_{\perp}), \|x\|_2 = 1} x^\top A x \\
        & =
        \max_{y \in \realR^{n-k}, \|y\|_2 = 1} (Q_{\perp} y)^\top A (Q_{\perp} y) \\
        & = \max_{y \in \realR^{n-k}, \|y\|_2 = 1}
        \left[
        \left(\epsilon \widetilde{M} y\right)^\top \Lambda_1
        \left(\epsilon \widetilde{M} y\right) \right. \\
        & \qquad\qquad\qquad\left.
        + \left(\left(I_{n-k} + \epsilon^2 \widetilde{N}\right) y\right)^\top
        \Lambda_2
        \left(\left(I_{n-k} + \epsilon^2 \widetilde{N}\right) y\right)
        \right] \\
        & \leq \lambda_{1} \left\|\epsilon \widetilde{M} y\right\|_2^2 + \lambda_{k+1} \left\|\left(I_{n-k} + \epsilon^2 \widetilde{N}\right) y\right\|_2^2\\
        & = \lambda_{1} \left\|\epsilon \widetilde{M} y\right\|_2^2 + \lambda_{k+1}\left(1 - \left\|\epsilon \widetilde{M} y\right\|_2^2\right)\\
        & \leq \lambda_{k+1} + \epsilon^2 (\lambda_1 -\lambda_{k+1}). 
    \end{aligned}
\]
Thus, we have
\[
    \begin{aligned}
        & \left| \lambda_{i}(Q_{F}^\top A Q_{F}) - \lambda_{i}(Q_{F}^\top A \langle Q \rangle Q_{F}) \right| \\
        & \leq
        \|F\|_2 \left(\frac{\|E\|_2}{\min|\lambda_i - \lambda(Q_{\perp}^\top A Q_{\perp})| - 2 \|F\|_2}\right)^2 \\
        & \leq
        \left(\frac{\|E\|_2}{\lambda_i - 3\lambda_{k+1} - 3 \epsilon^2 (\lambda_1 - \lambda_{k+1})}\right)^2 \|F\|_2 \\
        & \leq
        \left(\frac{\lambda_1 (1+ \epsilon^2) + \lambda_{k+1} \left(1 + \epsilon^2 \left(\frac{1}{1 + \sqrt{1-\epsilon^2}}\right)\right)}{\alpha_{i} \lambda_i}\right)^2 \left(\epsilon^2\lambda_{k+1} + \epsilon^4 (\lambda_1 - \lambda_{k+1}   )\right) \\
        & = C_{Nys,i} \cdot \left(\epsilon^2\lambda_{k+1} + \epsilon^4 (\lambda_1 - \lambda_{k+1})\right) \\
        & i = 1, ..., k,
    \end{aligned}  
\]
where $C_{Nys, i} := \left(\lambda_1 (1+ \epsilon^2) + \lambda_{k+1} \left(1 + \epsilon^2 \left(1 / 1 + \sqrt{1-\epsilon^2}\right)\right) / \alpha_{i} \lambda_i\right)^2$ and we assume $\lambda_i - 3\lambda_{k+1} - 3 \epsilon^2 (\lambda_1 - \lambda_{k+1}) = \alpha_i \lambda_i)$ for some $0 < \alpha_{i} < 1$ in the third inequality. This completes the proof.
\end{proof}
\begin{remark}
    \label{remark:Higher-order-of-Nystrom-SVD-vs-RR}
    For the leading few eigenvalues, i.e., all the $i$-th eigenvalues such that $\lambda_i \approx \lambda_1$, the \svdmethod's approximations usually provide about half of the RR error in the sense that
    \[
    |\lambda_i(A) - \sigma_i(AQ)| \lesssim \frac{1}{2} |\lambda_i(A) - \lambda_I(Q^\top A Q)|.
    \]
    This is due to
    \[
    \begin{aligned}
        |\lambda_i(A)^2 - \sigma_i(AQ)^2| & = |\lambda_i(A)^2 - \lambda_i(Q^\top A^2 Q)| \\ 
        & \leq \|\Lambda_1^2 - Q^\top A^2 Q\|_2 \\
        & = \|\epsilon^2 (M^\top \Lambda_1^2 + \Lambda_1^2 M + N^\top \Lambda_2^2 N) + \epsilon^4 M^\top \Lambda_1^2 M\|_2 \\
        & \leq \epsilon^2 (2 \lambda_1^2 + \lambda_{k+1}^2) + \epsilon^4 \lambda_1^2 
    \end{aligned}
    \]
    and
    \[
    \sigma_i(AQ) \geq \lambda_i(A) - C_{RR} \cdot \epsilon^2 \lambda_1,
    \]
    we have
    \begin{equation}
        \label{eq:Upper-bound-SVD}
        \begin{aligned}
            |\lambda_i(A) - \sigma_i(AQ)| & \leq \frac{|\lambda_i(A)^2 - \sigma_i(AQ)^2|}{\lambda_i(A) + \sigma_i(AQ)} \\ 
            & \leq \frac{\epsilon^2 (2 \lambda_1^2 + \lambda_{k+1}^2) + \epsilon^4 \lambda_1^2}{2\lambda_i(A) - C_{RR} \cdot \epsilon^2 \lambda_1}  \\
            & \leq \frac{\epsilon^2 (2 \lambda_1^2 + \lambda_{k+1}^2) + \epsilon^4 \lambda_1^2}{2\left(\lambda_i(A) - 2 \epsilon^2 \lambda_1\right)} \\
            & = \epsilon^2 \lambda_1 \left(2 + \epsilon^2 + \frac{\lambda_{k+1}^2}{\lambda_1^2}\right) \frac{\lambda_1}{2\left(\lambda_i(A) - 2 \epsilon^2 \lambda_1\right)} \\
            & \leq C_{RR} \cdot \epsilon^2 \lambda_1 \frac{\lambda_1}{2\left(\lambda_i(A) - 2 \epsilon^2 \lambda_1\right)} \\
            & = C_{SVD,i} \cdot \epsilon^2 \lambda_1.
        \end{aligned}
    \end{equation}
\end{remark}
\begin{remark}
    \label{remark:Higher-order-of-Nystrom}
    If $\lambda_1 / \lambda_i$, $\alpha_i$ are $O(1)$ and $\epsilon \ll 1$, then $C_{Nys}$ is $O(1)$. The first condition requires
    \begin{equation}
        \label{eq:Nys-higher-order-condition-1}
        C_i \lambda_i = \lambda_1, \quad \text{with moderate $O(1)$ constant } C_i
    \end{equation}
    and the second requires
    \begin{equation}
        \label{eq:Nys-higher-order-condition-2}
        \lambda_i \gg \lambda_{k+1}.
    \end{equation}
    Given the above conditions, we have
    \[
        C_{Nys,i} \leq \left(\frac{1+\epsilon^2}{\alpha_i}\right)^2 \left(C_i + \frac{1-\alpha_i}{3}\right)^2 \leq \left(\frac{2}{\alpha_i}\right)^2 \left(C_i + \frac{1-\alpha_i}{3}\right)^2.
    \]
    These imply that for some small $\tilde{k} < k$ such that conditions \eqref{eq:Nys-higher-order-condition-1} and \eqref{eq:Nys-higher-order-condition-2} are satisfied, the \Nystrom\ approximations to the eigenvalues $\lambda_1, ..., \lambda_{\tilde{k}}$ approximately $\lambda_{1}/\lambda_{k+1}$ times more accurate than the \svdmethod\ and RR. This can happen when the eigenvalues of $A$ decay rapidly.
\end{remark}



\subsection{Rayleigh-Ritz is better for trailing eigenvalue approximation}
Let us now discuss the accuracy of the methods for approximating the trailing eigenvalues. We first state a lemma that provides a lower bound for the RR approximations to the trailing eigenvalues.
\begin{lemma}
    \label{lemma:exact-vs-apprximate-eigenvalue-small}
    Suppose $Q$ is an orthonormal basis of $\subspace$, then
    \begin{equation}
        \lambda_i(Q^\top A Q) \geq \lambda_{n - \dim(\subspace) + i}(A), \quad (i = 1, ..., \dim(\subspace)).
    \end{equation}
\end{lemma}
If $\subspace$ is an approximation to the invariant subspace $\trailingsubspace$, then we have
\begin{equation}
    \lambda_i(A \langle Q \rangle) \geq \sigma_i(AQ) \geq \lambda_i(Q^\top A Q) \geq \lambda_{n - \dim(\subspace) + i}(A),
\end{equation}
where $i = 1, ..., \dim(\subspace)$. Thus, the trailing eigenvalue approximations 
\begin{align}
\hat{\lambda}_{j}^{\text{\rm(RR)}} &:= \lambda_{j - n + \dim(\subspace)}(Q^\top A Q), \nonumber \\
\hat{\lambda}_{j}^{\text{\rm(SVD)}} &:= \sigma_{j - n + \dim(\subspace)}(A Q), \nonumber \\
\hat{\lambda}_{j}^{\text{\rm(Nys)}} &:= \lambda_{j - n + \dim(\subspace)}\bigl(A \langle Q \rangle\bigr) \nonumber
\end{align}
have the reverse accuracy pattern
\begin{equation}
    \label{eq:smallest-eigenvalue-accuracy-relation}
    \left| \lambda_{j}(A) - \hat{\lambda}_{j}^{\text{(RR)}}\right| 
    \leq
    \left| \lambda_{j}(A) - \hat{\lambda}_{j}^{\text{(SVD)}}\right|
    \leq
    \left| \lambda_{j}(A) - \hat{\lambda}_{j}^{\text{(Nys)}}\right|,
\end{equation}
where $j = n-\dim(\subspace)+1, ..., n$. This shows that \Nystrom\  is not recommended when the target eigenvalues are the smallest ones. We present a remedy to this issue in Section~\ref{sec:trailing}. 

\subsection{Eigenvector approximation accuracy}
\label{subsec:Eigenvector-analysis-subsection}
We are currently unable to prove results that compare the accuracy of approximate eigenvectors from RR, SVD-$Q\tilde{V}$, SVD-$\tilde{U}$, and \Nystrom. Nonetheless, it is worth recalling that RR and SVD-$Q\tilde{V}$ seek eigenvectors from $\image(Q)$ while SVD-$\tilde{U}$ and \Nystrom\ seek from $\image(AQ)$. We can show that $\image(AQ)$ is a better subspace to extract eigenvectors than $\image(Q)$ when the dominant eigenspace is desired, whereas $\image(Q)$ is better when the trailing subspace is desired. In other words, we will illustrate that one extra subspace iteration $Q \mapsto AQ$ benefits the approximation to the invariant subspace $\leadingsubspace$ but harms that to the $\trailingsubspace$.

Let any $v \in \mathbb{R}^n$ be decomposed into $v = v_1 + v_2$, where $v_1 \in \mathcal{U}_1, v_2 \in \mathcal{U}_{1,\perp}$. We have $\projector_{\leadingsubspace} Av = A v_1 \in \mathcal{U}_1, \projector_{\mathcal{U}_{1,\perp}} Av = Av_2 \in \mathcal{U}_{1,\perp}$ since $\mathcal{U}_1$ is an invariant subspace of $A$, then
\[
    \begin{aligned}
        \sin^2 \angle(Av, \leadingsubspace) = \frac{\|\projector_{\mathcal{U}_{1,\perp}} Av\|_2^2}{\|Av\|_2^2} & = \frac{\|Av_2\|_2^2}{\|Av_1\|_2^2 + \|A v_2\|_2^2} \\
        & \leq \frac{\lambda_{k+1}^2 \|v_2\|_2^2}{\lambda_{k}^2 \|v_1\|_2^2 + \lambda_{k+1}^2 \|v_2\|_2^2} \\
        & = \frac{\|v_2\|_2^2}{\lambda_{k}^2 / \lambda_{k+1}^2 \|v_1\|_2^2 + \|v_2\|_2^2} \\
        & \leq \frac{\|v_2\|_2^2}{\|v_1\|_2^2 + \|v_2\|_2^2} = \sin^2 \angle(v, \leadingsubspace).
    \end{aligned}
\]
Where in the first inequality we used the facts that $\lambda_{k}^2 \|v_1\|_2^2 \leq \|A v_1\|_2^2$, $\|A v_2\|_2^2 \leq \lambda_{k+1}^2 \|v_2\|_2^2$ and $f(x) = x/(a + x)$ is increasing in $x > -a$. So,
\[
    \begin{aligned}
        \sin \angle(\image(AQ), \leadingsubspace) & = \sup_{v \in \image(Q), Av \neq 0} \sin \angle(Av, \leadingsubspace) \\
        & \leq \sup_{v \in \image(Q), v \neq 0} \sin \angle(v, \leadingsubspace) = \sin \angle(\image(Q), \leadingsubspace).
    \end{aligned}
\]
Similarly, for any $w \in \mathbb{R}^n$, we can decompose it into $w = w_1 + w_2$ with $w_1 \in \trailingsubspace, w_2 \in \mathcal{U}_{2,\perp}$ to show $\sin^2 \angle(Aw, \trailingsubspace) \geq \sin^2 \angle(w, \trailingsubspace)$ and further 
\[
    \sin \angle(\image(AQ), \trailingsubspace) \geq \sin \angle(\image(Q), \trailingsubspace).
\]


\section{Numerical experiments}
\label{sec:Numerical-illustrations-section}
Here we conduct experiments to illustrate the theoretical findings.
\subsection{Approximating the leading eigenpairs} 
Firstly, we illustrate Theorem~\ref{thm:Nystom-vs-SVD-vs-RR-Theorem}, and the leading eigenvalue approximation accuracy of the PSD matrix. For convenience, we scale the matrix such that $\|A\|_2=1$. In Figure~\ref{fig:leading-eigenvalue-accuracy-numerical-results-fast-decaying-spectrum} and Figure~\ref{fig:leading-eigenvalue-accuracy-numerical-results-slow-decaying-spectrum}, the spectrum of matrix $A$ decays exponentially and algebraically to $10^{-20}$, respectively. The orthogonal basis $Q$ is generated by the randomized rangefinder without power iteration \cite[Algorithm~9]{martinsson2020randomized} to approximate the invariant subspace $\leadingsubspace$ spanned by the leading eigenvectors of $A$. The right plot shows the approximate eigenvalues (Theorem~\ref{thm:Nystom-vs-SVD-vs-RR-Theorem}). The left plot illustrates the leading eigenvalue approximation accuracy, where the \Nystrom\ method performs the best, followed by \svdmethod, and RR is the least accurate, as predicted in Section~\ref{subsec:Leading-eigenvalue-subsection}. Also, we can observe a larger slope from the \Nystrom\ approximation error compared to the others when $A$'s spectrum is rapidly decaying. This phenomenon is less visible when the spectrum does not decay rapidly. 

Figure~\ref{fig:higher-order-of-Nystrom-figure} illustrates the upper bounds \eqref{eq:Upper-bound-RR}, \eqref{eq:Upper-bound-Nys}, and \eqref{eq:Upper-bound-SVD} for the eigenvalue approximation errors. It also demonstrates the higher-order accuracy phenomenon of \Nystrom\ for leading eigenvalue approximations: when the spectrum of $A$ decays rapidly, the leading few \Nystrom\ approximations can be more accurate than those produced by the \svdmethod\ and RR by a factor on the order of $\lambda_{1}(A)/\lambda_{k+1}(A)$; see Theorem~\ref{thm:Higher-order-accuracy-Nystrom-theorem}, Remark~\ref{remark:Higher-order-of-Nystrom-SVD-vs-RR} and Remark~\ref{remark:Higher-order-of-Nystrom}. Note that in Figure~\ref{fig:higher-order-of-Nystrom-figure} the orthonormal basis $Q$ is computed using \eqref{eq:transformed-Q-block-form}, rather than a randomized range finder. This allows us to control the subspace approximation error $\|\sin \Theta(\image(U_1), \image(Q))\|_2$ directly. In contrast, although a randomized range finder can capture the first few dominant eigenvectors accurately, the $k$th eigenvector is usually not, and thus the $\epsilon$ would be large, close to $1$, making the bounds less informative. 

\begin{figure}[htbp]
    \centering
    \includegraphics[scale = 0.3]{eigvals_leading_fast.eps}
    \caption{Accuracy of the leading eigenvalue approximations for \Nystrom, \svdmethod, and RR, where $n=1000$, $k=200$ and $A$ has an exponentially decaying spectrum from $1$ to $10^{-20}$. The figures illustrate Theorem~\ref{thm:Nystom-vs-SVD-vs-RR-Theorem} (right) and the approximate leading eigenvalues accuracy analysis (\ref{eq:largest-eigenvalue-accuracy-relation}) (left). A larger slope of the \Nystrom\ approximation error than others illustrates its higher-order accuracy with a fast-decaying spectrum.}
    \label{fig:leading-eigenvalue-accuracy-numerical-results-fast-decaying-spectrum}
\end{figure}

\begin{figure}[htbp]
    \centering
    \includegraphics[scale = 0.3]{Upper_bounds_new.eps}
    \caption{Numerical illustration for Theorem~\ref{thm:Higher-order-accuracy-Nystrom-theorem}, Remark~\ref{remark:Higher-order-of-Nystrom-SVD-vs-RR} and Remark~\ref{remark:Higher-order-of-Nystrom}. Only the leading few \Nystrom\ eigenvalue approximations are approximately $\lambda_{1}(A) / \lambda_{k+1}(A)$ (up to constants) more accurate than the \svdmethod\ and RR when $A$'s spectrum decays rapidly, where $\epsilon=0.01$, $n=1000$, $k=200$, $A$ has an exponentially decaying spectrum from $1$ to $10^{-20}$ and $Q$ is generated by \eqref{eq:eps-approximation-Q-formula}. The RR and \Nystrom\ bounds follow \eqref{eq:Upper-bound-RR}, \eqref{eq:Upper-bound-Nys} and the SVD bound here is $\min\left(C_{SVD,i} \cdot \epsilon^2 \lambda_1, C_{RR} \cdot \epsilon^2 \lambda_1\right)$. Note that the SVD and \Nystrom\ bounds depend on $i$, and \Nystrom\ bound exists if $\alpha_i > 0$.
    }
    \label{fig:higher-order-of-Nystrom-figure}
\end{figure}

\begin{figure}[htbp]
    \centering
    \includegraphics[scale = 0.3]{eigvals_leading_slow.eps}
    \caption{$A$ has a linearly decaying spectrum from $1$ to $10^{-20}$. The higher-order accuracy of the \Nystrom\ method is not seen.}
    \label{fig:leading-eigenvalue-accuracy-numerical-results-slow-decaying-spectrum}
\end{figure}

Next, we conduct numerical experiments for the leading eigenvector approximation accuracy. The accuracy of eigenvector approximations are measured by $\sin \angle (u_{i}, \hat{u}_{i})$. Figure~\ref{fig:leading-eigenvector-accuracy-numerical-results} illustrates the accuracy when $A$ has an exponentially and an algebraically decaying spectrum, respectively. We can see a similar accuracy pattern to that in the leading eigenvalue approximation if $A$ has a fast-decaying spectrum. In contrast, none of the methods provides accurate approximations if $A$'s spectrum does not decay rapidly; this is because then the randomized rangefinder captures the invariant subspace $\leadingsubspace$ less accurately. Also, \Nystrom\ and SVD-$\tilde{U}$ provide better approximations than SVD-$Q\tilde{V}$ and RR, illustrating that $\image(AQ)$ is a better choice for finding the approximate leading eigenvectors than the given subspace $\image(Q)$, as we discused in Section~\ref{subsec:Eigenvector-analysis-subsection}.

\begin{figure}[htbp]
    \centering
    \includegraphics[scale = 0.3]{eigvecs_leading.eps}
    \caption{Accuracy of the leading eigenvector approximations, $k=200$ and $A$ has an exponentially (left) and an algebraically (right) decaying spectrum from $1$ to $10^{-20}$. The SVD-$\tilde{U}$ and \Nystrom\ methods provide more accurate approximations than the others, illustrating that $\image(AQ)$ is a better subspace to extract the leading eigenvectors.}
    \label{fig:leading-eigenvector-accuracy-numerical-results}
\end{figure}

\subsection{Approximating trailing eigenpairs}\label{sec:trailing}
Finally, we numerically examine the trailing eigenvalue approximation accuracy analysis (Eq.~\eqref{eq:smallest-eigenvalue-accuracy-relation}) and the performance in trailing eigenvector approximation. We use the same settings as in the experiments for the leading one, except for the spectrum-decay speed. We set the spectrum to decay algebraically and linearly to $10^{-20}$ rather than exponentially. The orthonormal basis $Q$ is generated by perturbing the exact orthonormal basis of the invariant subspace $\trailingsubspace$, i.e., $U_2 + \epsilon G$, where $G$ is a Gaussian matrix with i.i.d entries of zero mean and $1/n$ variance, i.e., $G_{ij} \sim \mathcal{N}(0, 1/n)$. Figure~\ref{fig:trailing-eigenvalue-accuracy-numerical-results} shows the reverse accuracy pattern (Eq.~\eqref{eq:smallest-eigenvalue-accuracy-relation}) in trailing eigenvalue approximation. Figure~\ref{fig:trailing-eigenvector-accuracy-numerical-results} illustrates the accuracy of the trailing eigenvector approximation, where only RR can provide good approximations since $\image(Q)$ can better capture $\trailingsubspace$ than $\image(AQ)$.

\paragraph{Remedy for \Nystrom\ when approximating trailing eigenvalues}
Instead of directly approximating the trailing eigenpairs of $A$, a possible remedy is to exploit the superiority of $\image(AQ)$ in approximating $\leadingsubspace$ (compared to $\image(Q)$) by computing the leading eigenpairs of the shifted and negated matrix $-A + \gamma I$, where $\gamma \geq \lambda_1(A)$, where $\gamma$ can be estimated for example using a small number of iterations of the Lanczos method, together with Weyl's bound using the residual: $\hat{\lambda}_1+\|A\hat{u}_1-\hat{\lambda}_1\hat{u}_1\|_2$ is a practical estimate for the upper bound. Figure~\ref{fig:trailing-eigenvalue-accuracy-shift-version-numerical-results} and Figure~\ref{fig:trailing-eigenvector-accuracy-shift-version-numerical-results} illustrate the accuracy of the eigenvalue and eigenvector approximations using the shift trick, respectively. \Nystrom's method provides better approximations based on RR, and thus better accuracy than the case of the unshifted version.

\begin{figure}[htbp]
    \centering
    \includegraphics[scale = 0.3]{eigvals_trailing_new.eps}
    \caption{Accuracy comparison of the trailing eigenvalue approximations, where $n=1000$, $k=200$, and $A$ has an algebraically decaying spectrum (left) and a linearly decaying spectrum (right) from $1$ to $10^{-20}$. A reverse accuracy pattern is observed compared to the leading eigenvalue approximation case.}
    \label{fig:trailing-eigenvalue-accuracy-numerical-results}
\end{figure}

\begin{figure}[htbp]
    \centering
    \includegraphics[scale = 0.3]{eigvec_trailing_new.eps}
    \caption{Accuracy comparison of the trailing eigenvector approximations, where $k=200$ and $A$ has an algebraically decaying spectrum (left) and a linearly decaying spectrum (right) from $1$ to $10^{-20}$. Only RR provides good approximations.}
    \label{fig:trailing-eigenvector-accuracy-numerical-results}
\end{figure}

\begin{figure}[htbp]
    \centering
    \includegraphics[scale = 0.3]{eigval_trailing_shift_new.eps}
    \caption{Accuracy of the trailing eigenvector approximation using the shift trick to remedy, where $k=200$ and $A$ has an algebraically decaying spectrum (left) and a linearly decaying spectrum (right) from $1$ to $10^{-20}$. The leading eigenvalues of $-A + \gamma I$ are clustered around $\gamma$. We can benefit from the \Nystrom\ by applying it to the shifted matrix.}
    \label{fig:trailing-eigenvalue-accuracy-shift-version-numerical-results}
\end{figure}

\begin{figure}[htbp]
    \centering
    \includegraphics[scale = 0.3]{eigvec_trailing_shift_new.eps}
    \caption{Accuracy of the trailing eigenvector approximations using the shift trick to remedy.}
    \label{fig:trailing-eigenvector-accuracy-shift-version-numerical-results}
\end{figure}


\section{Conclusion and future work}
We have shown how to exploit the PSD property in the extraction step of the subspace methods to improve the accuracy of leading eigenvalue approximations beyond RR via \svdmethod\ and especially \Nystrom. Further, the accuracy of \Nystrom\ is approximately $\lambda_1(A) / \lambda_{k+1}(A)$ times better, where $k$ is the dimensional of the given subspace, compared to the RR and \svdmethod\ for the leading $\tilde{k}$ eigenvalues $(\tilde{k} < k)$ if the matrix satisfies additional conditions, for example, a fast-decaying spectrum. The situation is the opposite for trailing eigenpair approximation. However, we can fix this by working with the matrix $-A + \gamma I$ with $\gamma \geq \lambda_1(A)$. 

Several directions remain for future investigation. First, it would be useful to extend the present analysis of the higher-order accuracy of \Nystrom\ to alternative metrics of subspace approximation accuracy. For example, the projection residual metric $\|(I- QQ^\top)A\|$ that captures the matrix approximation quality rather than only the worst-direction subspace error, as in Section~\ref{subsec:Analysis-High-order-accuracy-Nystrom}, may lead to sharper and more informative eigenvalue error estimates. Second, the current analysis focuses on exterior eigenvalues. A natural next step is to investigate how additional matrix structure, including but not limited to positive semi-definiteness, can be exploited to improve the computation of interior eigenvalues. Third, it would be interesting to study whether analogous accuracy improvements can be obtained for generalized Rayleigh--Ritz approximations in generalized eigenvalue problems.

\bibliographystyle{siamplain}
\bibliography{references}
\end{document}

\typeout{get arXiv to do 4 passes: Label(s) may have changed. Rerun}
